\documentclass[10pt]{amsart}
\usepackage[a4paper,margin=23mm]{geometry}
\usepackage{amsmath,amssymb,amsthm,mathtools,hyperref,graphicx,xspace,letltxmacro,xpatch}
\usepackage[english]{babel}
\newtheorem{theo}{Theorem}[section]
\newtheorem{lemm}[theo]{Lemma}\newtheorem{prop}[theo]{Proposition}\newtheorem{coro}[theo]{Corollary}\newtheorem{conj}[theo]{Conjecture}
\theoremstyle{definition}\newtheorem{defi}[theo]{Definition}\newtheorem{rema}[theo]{Remark}\newtheorem{exam}[theo]{Example}\newtheorem{exem}[theo]{Example}\newtheorem{nota}[theo]{Notation}
\newcommand{\ie}{i.e.}\newcommand{\eg}{e.g.}\newcommand{\resp}{resp.}\newcommand{\skpt}{\par\smallskip\noindent}
\emergencystretch=3em\pagestyle{plain}
\usepackage{subfig}
\usepackage{xpatch}

\DeclareMathOperator{\E}{\mathbf{E}}



\newcommand{\bbN}{\mathbb{N}}
\newcommand{\Z}{\mathbb{Z}}
%\newcommand{\R}{\mathbb{R}}
%\newcommand{\C}{\mathbb{C}}
\newcommand{\T}{\mathbb{T}}
\newcommand{\ind}[1]{\mathbf{1}_{\left\{#1\right\}}}
%\newcommand{\indset}[1]{\mathbf{1}_{#1}}
\newcommand{\floor}[1]{{\left\lfloor #1 \right\rfloor}}
%\newcommand{\ceil}[1]{{\left\lceil #1 \right\rceil}}

\renewcommand{\P}{\mathbf{P}}


\newcommand{\defeq}{:=}

\renewcommand{\S}{\mathcal{S}}
\newcommand{\clH}{\mathcal{H}}
\newcommand{\clF}{\mathcal{F}}
\newcommand{\clG}{\mathcal{G}}
%%%%%%%%%%%%%%%%%%%%%%%%%%%%%%%%%%%%%%%%%%%%%%%%%%%%%%%%%%%
%~ This fixes spacing issues around \left( and \right) and other similar symbols
\let\originalleft\left 
\let\originalright\right
\renewcommand{\left}{\mathopen{}\mathclose\bgroup\originalleft}
\renewcommand{\right}{\aftergroup\egroup\originalright}
%%%%%%%%%%%%%%%%%%%%%%%%%%%%%%
\let\oldin\in
\renewcommand{\in}{\mathchoice{\oldin}{\oldin}{\,\oldin\,}{\,\oldin\,}}
%%%%%%%%%%%%%%%%%Arrow
\renewcommand*{\to}{\mathchoice{\longrightarrow}{\rightarrow}{\rightarrow}{\rightarrow}}
\let\oldmapsto\mapsto
\renewcommand*{\mapsto}{\mathchoice{\longmapsto}{\oldmapsto}{\oldmapsto}{\oldmapsto}}

%%%%%%%%%%%%%%%%%%%%%%%%%%%%%%%%%%%%%%%%%%%%%%%%%%%%%%%%%%%%%%%%%%%%
\graphicspath{{./figures/}}

\newcommand*{\mk}{\mkern -1mu}
\newcommand*{\Mk}{\mkern -2mu}
\newcommand*{\mK}{\mkern 1mu}
\newcommand*{\MK}{\mkern 2mu}

%\hypersetup{urlcolor=purple, linkcolor=blue, citecolor=red}

\newcommand*{\relabel}{\renewcommand{\labelenumi}{(\theenumi)}}
\newcommand*{\romanenumi}{\renewcommand*{\theenumi}{\roman{enumi}}\relabel}
\newcommand*{\Romanenumi}{\renewcommand*{\theenumi}{\Roman{enumi}}\relabel}
\newcommand*{\alphenumi}{\renewcommand*{\theenumi}{\alph{enumi}}\relabel}
\newcommand*{\Alphenumi}{\renewcommand*{\theenumi}{\Alph{enumi}}\relabel}
\let\oldtilde\tilde
\renewcommand*{\tilde}[1]{\mathchoice{\widetilde{#1}}{\widetilde{#1}}{\oldtilde{#1}}{\oldtilde{#1}}}
\let\oldhat\hat
\renewcommand*{\hat}[1]{\mathchoice{\widehat{#1}}{\widehat{#1}}{\oldhat{#1}}{\oldhat{#1}}}

\newcommand{\llbracket}{\mathopen{[\mkern -2.7mu[}}
\newcommand{\rrbracket}{\mathclose{]\mkern -2.7mu]}}
\newcommand{\llbracketclose}{\mathclose{[\mkern -2.7mu[}}
\newcommand{\rrbracketopen}{\mathopen{]\mkern -2.7mu]}}

\newcommand{\llb}{\llbracket}
\newcommand{\rrb}{\rrbracket}
\newcommand{\llbigbracket}{\mathopen{\bigl[\mkern -4.2mu\bigl[}}
\newcommand{\rrbigbracket}{\mathclose{\bigr]\mkern -4.2mu\bigr]}}
\newcommand{\llbg}{\llbigbracket}
\newcommand{\rrbg}{\rrbigbracket}
%%%%%%%%%%%%%%%%%%%%%%%%%%%%%%%%%%%%%%%%%%%%%%%%%%%%%%%%%%%%%%%%%%%%

\makeatletter
\expandafter\def\csname externalRef@fig:IBM\endcsname{2.1}
\expandafter\def\csname externalRef@fig:couplage\endcsname{3.1}
\expandafter\def\csname externalRef@thm:finitevinfinite\endcsname{1.9}
\expandafter\def\csname externalRef@theo1.9.1\endcsname{1}
\LetLtxMacro\NativeRef\ref
\newcommand{\SourceRef}[1]{\ifcsname externalRef@#1\endcsname\href{https://ahl.centre-mersenne.org/item/AHL_2025__8__635_0/}{\csname externalRef@#1\endcsname\ of the source article}\else\NativeRef{#1}\fi}
\AtBeginDocument{\LetLtxMacro\NativeRef\ref\LetLtxMacro\ref\SourceRef}
\makeatother
\begin{document}
\title{A zero-one freezing law for the infinite-bin model with non-increasing reproduction probabilities}
\author{Bastien Mallein \and Sanjay Ramassamy \and Arvind Singh}
\maketitle
\section{Introduction}\label{sec:introduction}

The infinite-bin model (IBM) is a discrete-time particle system on $\Z$ introduced by Foss and Konstantopoulos~\cite{FK}. In this process, at each integer time, a particle chosen according to its rank reproduces by creating a newborn particle one step to its right. Informally, it is constructed as a balls and bins scheme, by placing a bin at each site of $z\in\Z$ with a certain number of balls, corresponding to the particles at initial time. We assume that all the bins at location $z$ large enough are initially empty, so it is possible to number all the balls in the system starting from the rightmost one. We define the dynamics of the infinite-bin model as follows. Let $\mu$ be a probability distribution on $\bbN
:= \{1,2,\ldots\}$. At each integer time, the $k^{\rm th}$ rightmost ball is selected with probability $\mu(k)$ and a ball is added to the bin immediately on its right. See Figure~\ref{fig:IBM} for an illustration of the evolution of the IBM.

We introduce some notation needed to define the infinite-bin model formally. The state space of the IBM is defined as
\[
\S
:= \bigl\{ X : \Z \to \Z_+\cup \{\infty\} \text{ such that } \sup\{ j \in \Z : X(j) \neq 0 \} \in (-\infty,\infty) \bigr\}.
\]
An element $X\in \S$ represents a \emph{configuration of balls} where, for each $k \in \Z$, the bin at position $k$ contains $X(k)$ balls.

Let us point out that the total number of balls in a configuration $X \in \S$ may be finite or infinite (but cannot be zero). We even allow the number of balls in a given bin $k$ to be infinite i.e. $ X(k) = \infty$. For the dynamics of the infinite bin model, such a bin will represent an impassable barrier (the configuration on the left of this bin will not matter and remain unchanged forever). A configuration $X \in \S$ is said to be
\begin{itemize}
\item \emph{locally finite} if $X(k) < \infty$ for all $k\in \Z$,
\item \emph{bounded} if $\sup_{k\in \Z} X(k) < \infty$.
\end{itemize}

Given $\mu$ a probability distribution on $\bbN$ and $X_0 \in \S$ an initial configuration, an infinite-bin model $X$ with reproduction law $\mu$ starting from $X_0$ is constructed as an $\S$-valued Markov process in the following fashion. Given $(\xi_n, n \in \bbN)$ i.i.d. random variables with law $\mu$, at each integer time $n$, $X_n$ is constructed from $X_{n-1}$ by adding a new ball to the bin immediately to the right of the bin containing the $\xi_{n}^{\rm th}$ rightmost ball. In other words, writing $J_{n-1}$ for the index of the bin containing the $\xi_{n}^{\rm th}$ rightmost ball at time $n-1$, defined by
\[
\sum_{j = J_{n-1}}^\infty X_{n-1}(j) \geq \xi_{n} > \sum_{j=J_{n-1} +1}^\infty X_{n-1}(j),
\]
we set\footnote{If such an index does not exist, which is the case if the configuration $X_{n-1}$ contains fewer than $\xi_{n}$ balls, we just set $X_{n}=X_{n-1}$.} $X_{n}(\cdot) = X_{n-1}(\cdot) + \ind{\cdot\,=\,J_{n-1}+1}$.

In contrast, we focus here on a different question and study instead how the number of balls in a given bin grows as time passes. Obviously, for all $k \in \Z$, the number $X_n(k)$ of balls in the bin at position $k$ at time $n$ is non-decreasing in $n$ (since balls are never removed from bins). Hence we can define the random variables
\[
X_\infty(k) = \lim_{n \to \infty}\uparrow X_n(k) \in \Z_+ \cup \{\infty\} \quad \text{a.s.}
\]
We call the event $\{X_\infty(k) < \infty\}$ the \emph{freezing} of bin $k$, as it means that, on this event, there is some finite (random) time after which no ball is ever added to bin $k$. This paper aims to describe conditions on the distribution $\mu$ and the initial configuration $X_0$ for which freezing occurs.
\setcounter{theo}{2}
\begin{theo}\label{thm:01law}
Assume that the probability distribution $\mu$ satisfies
\begin{equation}\label{eqn:monotonous}
\text{$(\mu(n), n \geq 1)$ is non-increasing.}
\end{equation}
Then, for any initial configuration $X_0 \in \S$, we have
\[
\P(X_\infty(k) < \infty) \in \{0,1\} \quad \text{for all $k\in \Z$.}
\]
\end{theo}
\section{General properties of infinite-bin models}\label{sec:preliminaries}

Before studying properties of the infinite-bin model dynamics, we introduce a few extra notations that will allow us to describe in greater details the local dynamic of the model.

Given an infinite-bin model configuration $X \in \S$, we define its \emph{rightmost barrier} as
\begin{equation}\label{def:deltabarrier}
\delta(X)
:= \sup\{j\in \Z : X(j) = +\infty\} \;\in \Z\cup\{-\infty\}
\end{equation}
with the convention $\delta(X) = \sup \emptyset = -\infty$ when $X$ is locally finite. The \emph{front} of $X$ is the location of its rightmost non-empty bin (which is well-defined by definition of a configuration), i.e.
\[
F(X)
:= \max\{ j \in \Z : X(j) \neq 0\} \in \Z.
\]
For all $k \in \bbN$, we also define $B(X,k)$ as the location of the $k^{\rm th}$ rightmost ball in $X$, i.e.
\[
B(X,k)
:= \sup\left\{ j \in \Z : \sum_{i=j}^\infty X(i) \geq k \right\},
\]
again with $\sup \emptyset = -\infty$. We observe that
\[
F(X) = B(X,1)\quad\text{and}\quad\delta(X) = \lim_{k\to\infty} B(X,k).
\]

\begin{rema}
For all $X \in \S$ and $k \in \bbN$, the quantity $B(X, k)$ is well-defined irrespectively of the relative ordering of the balls inside each bin. In other words, the dynamic of the infinite-bin model does not require to fix an ordering of the balls. However, it will be convenient to specify an ordering for the balls in each bin, in order to introduce a genealogical structure to the particle system.
\end{rema}


To give a formal definition of the infinite-bin model dynamics, we introduce the operator $\Phi_\xi : \S \to \S$ that maps a configuration to the new configuration obtained by adding a single new ball immediately to the bin on the right of the $k^{\rm th}$ rightmost ball:
\begin{equation}\label{eq:operatorPhi}
\Phi_\xi(X)
:=
\begin{cases}
\left(X(j) + \ind{j\,=\,B(X,\xi)+1}, j \in \Z\right) & \text{ if } B(X,\xi) > -\infty,\\
X & \text{ otherwise.}
\end{cases}
\end{equation}
Given a probability distribution $\mu$ and an initial configuration $X_0 \in \S$, the infinite-bin model $(X_n,\, n\geq 0)$ with distribution $\mu$ starting from $X_0 \in \S$ can be constructed through the following recursion equation, for $n\geq 0$
\begin{equation}\label{eqn:ibm}
X_{n+1}
:= \Phi_{\xi_{n+1}} (X_n),
\end{equation}
where $(\xi_n, n \geq 1)$ is an i.i.d. sequence of random variables with law $\mu$.


In the rest of this article, we always picture balls as stacked vertically in their bin, and ranked from bottom to top. When a new ball is added, it is placed on the top of the previous balls in the bin. With this ordering, for each ball in $X$, we can speak about
\begin{itemize}
\item its \emph{rank} which corresponds to its index when all balls are ordered starting from the front, from right to left, and then from bottom to top inside each bin.
\item Its \emph{location} which is the index/position of the bin it belongs to.
\item Its \emph{height} which corresponds to its index inside the bin with the balls ordered from bottom to top.
\end{itemize}
Let us note that, by definition, for each $\xi \in \bbN$ the location and heights of the balls already present in $X$ are unchanged in the new configuration $\Phi_\xi(X)$. However, the ranks of the balls in $\Phi_\xi(X)$ are either unchanged or increased by one. See Figure~\ref{fig:IBM} for an illustration of the procedure.
\section{Coupling of IBMs for different initial configurations}\label{sec:coupling}

In this section, we prove Theorem~\ref{thm:01law} and Theorem~\ref{thm:finitevinfinite}$\MK$\eqref{theo1.9.1}. Both results make repetitive use of couplings of IBMs starting from different initial configurations. We start with the following easy lemma which relates the evolution of two IBMs under the usual coupling.

\begin{lemm}\label{lemma:couplage1}
Let $X_0, \tilde{X}_0 \in \S$ denote two initial configurations with the same position $\delta$ for their rightmost infinite barrier (if it exists):
\[
\delta
:= \delta(X_0) = \delta(\tilde{X}_0) \in \Z\cup \{-\infty\}.
\]
Fix $j_0 > \delta$ and assume that the following condition holds true at time $n=0$
\begin{equation}
\text{for all }j \in \rrbracketopen \delta, j_0 \rrb,\quad\tilde{X}_n(j) = X_n(j) \quad\text{and}\quad \sum_{k =j}^{\infty} X_n(k) = \sum_{k= j}^{\infty} \tilde{X}_n(k).\tag{H1}\label{hypH1}
\end{equation}
Let $(X_n)$ and $(\tilde{X}_n)$ be defined by the recursion equation~\eqref{eqn:ibm} using the same sequence $(\xi_n)$. Then~\eqref{hypH1} holds true for all times $n$. Furthermore, we have, for all $n\geq 0$ and all $j\in \rrbracketopen \delta, j_0 + 1\rrb$,
\begin{equation}\label{incH1}
\tilde{X}_{n+1} (j) = \tilde{X}_{n} (j) + 1 \quad \Longleftrightarrow \quad X_{n+1} (j) = X_{n} (j) + 1.\qquad\text{a.s.}
\end{equation}
\end{lemm}
\begin{proof}
We prove the result by induction on $n$. It is true for $n=0$ by assumption. Suppose now that~\eqref{hypH1} holds for a given $n\geq 0$ and consider $\xi_{n+1}$. There are two cases:
\begin{itemize}
\item If $\xi_{n+1} \leq \sum_{k =j_0}^{\infty} X_n(k) = \sum_{k= j_0}^{\infty} \tilde{X}_n(k)$, then we add a ball to $X_n$ and a ball in $\tilde{X}_n$, possibly at different locations, but always in a bin strictly on the right of $j_0$ in both cases. Therefore, \eqref{hypH1} holds for $n+1$ (as well as~\eqref{incH1} because no ball is added to bin $k$).
\item If $\xi_{n+1} > \sum_{k =j_0}^{\infty} X_n(k) = \sum_{k= j_0}^{\infty} \tilde{X}_n(k)$, then thanks to~\eqref{hypH1}, the ball ranked $\xi_{n+1}$ in $X_n$ is at the same location and same height as the ball ranked $\xi_{n+1}$ in~$\tilde{X}_n$. Thus, the ball added at time $n+1$ is at the same location in $X_{n+1}$ and $\tilde{X}_{n+1}$ hence~\eqref{hypH1} and~\eqref{incH1} still hold.\qedhere
\end{itemize}\let\qed\relax
\end{proof}

We can improve on the previous lemma by coupling in a larger class of initial configurations if we assume that the distribution $\mu$ is monotonic.


\begin{lemm}\label{lemma:couplage2}
Assume that $(\mu(n), n \geq 1)$ is non-increasing. Let $X_0, \tilde{X}_0 \in \S$ denote two initial configurations with the same position $\delta$ for their rightmost infinite barrier (if it exists):
\[
\delta
:= \delta(X_0) = \delta(\tilde{X}_0) \in \Z\cup \{-\infty\}.
\]
Fix $j_0 > \delta$ and assume that the following condition holds true at time $n=0$
\begin{equation}\label{hypH2}
\text{for all }j \in \rrbracketopen \delta, j_0 \rrb,\quad\tilde{X}_n(j) \leq X_n(j) \quad\text{and}\quad \sum_{k =j}^{\infty} X_n(k) \leq \sum_{k= j}^{\infty} \tilde{X}_n(k).\tag{H2}
\end{equation}
then, there exists a coupling of two infinite-bin models $(X_n)$ and $(\tilde{X}_n)$ with distribution $\mu$, starting respectively from $X_0$ and $\tilde{X}_0$ such that~\eqref{hypH2} holds for all times $n$. Furthermore, this coupling also satisfies that, for $n\geq 0$ and all $j \in \rrbracketopen \delta, j_0 + 1\rrb$,
\begin{equation}\label{incH2}
\tilde{X}_{n+1} (j) = \tilde{X}_{n} (j) + 1 \quad \Longrightarrow \quad X_{n+1} (j) = X_{n} (j) + 1.\quad\text{a.s.}
\end{equation}
\end{lemm}

See Figure~\ref{fig:couplage} for an illustration of two configurations $X$ and $\tilde{X}$ that satisfy assumption~\eqref{hypH2}.

\goodbreak
\begin{rema}
If $X_0$ and $\tilde{X}_0$ satisfy~\eqref{hypH2}, then the same holds true after applying either of the following two elementary operations:
\begin{enumerate}
\item\label{rema3.3.1} We add a ball in $\tilde{X}_0$ at a location $j > j_0$.
\item\label{rema3.3.2} We move a ball of $\tilde{X}_0$ from a location $j \leq j_0$ to a new location $j' > j_0$.
\end{enumerate}
\end{rema}

\begin{proof}
We just need to construct the coupling for the first step and then we can proceed by induction afterward thanks to the Markov property. Let
\[
N
:= \sum_{j > \delta} X_0(j) \quad\text{and} \quad \tilde{N}
:= \sum_{j > \delta} \tilde{X}_0(j)
\]
denote the total number of balls in configuration $\tilde{X}$ and $\tilde{X}_0$ respectively. We have $N\leq \tilde{N}\leq\infty$ according to~\eqref{hypH2} (with quantities being possibly infinite when $\delta(X) =-\infty$). We also define the number of balls strictly on the right of bin $j_0$:
\[
K
:= \sum_{j = j_0 + 1}^\infty X_0(j) \quad \text{and} \quad \tilde{K}
:= \sum_{j = j_0 + 1}^\infty \tilde{X}_0(j).
\]
We also have $K \leq \tilde{K} < \infty$ thanks to~\eqref{hypH2}. Furthermore, if $K = \tilde{K}$ then, we must have $\tilde{X}_0(j) = X_0(j)$ for all $j\in \rrbracketopen \delta, j_0\rrb$. This exactly means that~\eqref{hypH1} is satisfied so we are back in the setting of the previous lemma and we can use the trivial coupling to prove the result at the next step. Thus, we now assume that $K < \tilde{K}$.

We also assume without loss of generality that $\mu(\llb 1, K\rrb) < 1$ otherwise the result is again trivial. Because $\mu$ is non-increasing, this assumption is equivalent to requiring $\mu(K+1) > 0$.

\goodbreak
Let $\xi$, $\chi$ and $U$ denote three independent random variables such that
\begin{itemize}
\item $\xi$ is distributed according to $\mu$.
\item $\chi$ is distributed according to the conditional law $\mu(\cdot \mid \rrbracketopen K, \tilde{K}\rrb)$.
\item $U$ is uniform in $[0,1]$.
\end{itemize}
We construct a random variable $\tilde{\xi}$ depending on $\xi$ in the following way:
\begin{enumerate}\alphenumi
\item\label{casea} If $\xi \leq K$, we set $\tilde{\xi}
:= \xi$.
\item\label{caseb} If $\xi > \tilde{N}$, we set $\tilde{\xi}
:= \xi$.
\item\label{casec} If $\xi \in \rrb N, \tilde{N}\rrb$, we set $\tilde{\xi}
:= \chi$.
\item\label{cased} If $\xi = k \in \rrb K,N\rrb$, then $k$ is the rank of a ball in $X_0$. If a ball with the same location and same height also exists in configuration $\tilde{X}_0$, then we denote $\tilde{k}$ its rank in $\tilde{X}_0$. We set
\[
\tilde{\xi}
:=
\begin{cases}
\tilde{k} & \text{if }\tilde{k}\text{ exists and }U < \frac{\mu(\tilde{k})}{\mu(k)},\\
\chi & \text{otherwise.}
\end{cases}
\]
\end{enumerate}
We have $K \leq N \leq \tilde{N}$ so the four cases~\eqref{casea}, \eqref{caseb}, \eqref{casec} and~\eqref{cased} are disjoint and $\tilde{\xi}$ is well-defined. We now define the coupling:
\[
X_1
:= \Phi_\xi(X_0) \quad \text{and}\quad \tilde{X}_1
:= \Phi_{\tilde\xi}(\tilde{X}_0).
\]
where $\Phi$ is the operator defined in~\eqref{eq:operatorPhi}. We check that~\eqref{hypH2} holds for $n=1$. Indeed, we observe that
\begin{itemize}
\item In case~\eqref{casea}, configurations $X_1$ and $\tilde{X}_1$ are obtained by adding a ball into a bin, not necessarily the same, but located in $\llb j_0 + 2, +\infty\llbracketclose $ in both cases.
\item In case~\eqref{caseb}, we have $X_1 = X_0$ and $\tilde{X}_1 = \tilde{X}_0$ if $\delta = -\infty$. Otherwise, the new ball is added in both cases to the bin $\delta +1$.
\item In case~\eqref{casec}, we have $X_1 = X_0$ (if $\delta =-\infty$) or $X_1(\delta+1) = X_0(\delta+1)+1$\linebreak (if $\delta > -\infty$), while $\tilde{X}_1$ is obtained by adding a ball to $\tilde{X}_0$ into a bin located somewhere in $\llb j_0+2,\;+\infty\llb \subset \rrb \delta,\infty\llbracketclose$.
\item In case~\eqref{cased}, if $\tilde{k}$ exists and $U < \mu(\tilde{k}) / \mu(k)$ then $X_1$ and $\tilde{X}_1$ are obtained by adding a ball in the same bin for both configuration. Otherwise, a ball is added in a bin located somewhere in $\rrbracketopen \delta, j_0 + 1\rrb$ for configuration $X_0$ whereas a ball is added in a bin located somewhere in $\llb j_0 + 2, +\infty\llbracketclose$ for configuration~$\tilde{X}_0$.
\end{itemize}
In all cases, hypothesis~\eqref{hypH2} is preserved. It is also clear that~\eqref{incH2} holds because a ball is added in $\tilde{X}$ at some location $\leq j_0 + 1$ when $\tilde{\xi} \in \rrbracketopen\tilde{K}, \tilde{N}\rrb$ but this can only occur in case~\eqref{cased} and a ball is always added at the same location in $X$ when this happens.

It remains to prove that $\tilde{\xi}$ has distribution $\mu$.
\begin{itemize}
\item Let $i \leq K$ or $i > \tilde{N}$, we have by construction
\[
\P(\tilde{\xi} = i) = \P(\xi = i) = \mu(i).
\]
\item Let $i \in \rrbracketopen\tilde{K}, \tilde{N}\rrb$. We observe that the ball with rank $i$ in $\tilde{X}_0$ is in a bin at some location $\rrbracketopen\delta, j_0\rrb $. Because of Assumption~\eqref{hypH2}, a ball with the same location{\pagebreak\linebreak\ }
and same height also exists in $X_0$ and its corresponding rank $k$ in $X_0$ satisfies $k \leq i$. Therefore, we get
\[
\P(\tilde{\xi} = i) = \P\left(\xi = k,\; U < \frac{\mu(i)}{\mu(k)}\right) = \mu(k) \P\left(U < \frac{\mu(i)}{\mu(k)}\right) = \mu(i),
\]
where we used the fact $0 \leq \mu(i) / \mu(k) \leq 1$ because $(\mu(n))$ is non-increasing.
\item Let $i \in \rrbracketopen K, \tilde{K}\rrb$. We can only have $\tilde{\xi} = i$ in cases~\eqref{casec} or~\eqref{cased} when $\tilde{\xi} = \chi$. But by construction, the event $\{\tilde{\xi} = \chi\}$ depends only on $\xi$ and $U$ which are independent of $\chi$. Therefore,
\[
\P\left(\tilde{\xi} = i\right) = \P\left(\tilde{\xi} = \chi,\; \chi = i\right) = \P\left(\tilde{\xi} = \chi\right) \P(\chi = i) = \alpha \frac{\mu(i)}{\mu\left(\rrbracketopen K, \tilde{K}\rrb\right)}
\]
for some constant $\alpha \in [0,1]$. Finally, since $\tilde{\xi}$ is a non-defective random variable, we have $\sum \P(\tilde{\xi} = i) = 1$ which implies that $\alpha = \mu(\rrbracketopen K, \tilde{K}\rrb)$ and the proof of the Lemma~\ref{lemma:couplage2} is complete. \qedhere
\end{itemize}
\end{proof}


\subsection{\texorpdfstring{$0-1$}{0-1} law for the freezing of a bin}

We assume here that $(\mu(n), n \geq 1)$ is non-increasing and we show that, for any initial configuration $X_0 \in \S$ and any $k\in \Z$, we have
\begin{equation*}
\P(X_\infty(k) = \infty) \in \{0, 1\}.
\end{equation*}
\begin{proof}[Proof of Theorem~\ref{thm:01law}]
Let $k \in \bbN$, for all $n \in \bbN$, we set
\[
B_n
:= \left\{ X_{n}(k) > X_{n-1}(k)\right\},
\]
the event when the $n^{\rm th}$ ball added in the process goes into bin $k$. We prove that the events $(B_n, n \geq 1)$ are negatively correlated, which is enough to complete the proof of the $0-1$ law. Indeed, observe that if
\[
\sum_{n \geq 1} \P(B_n) < \infty,
\]
then by Borel--Cantelli lemma, we have $X_n(k) = X_{n-1}(k)$ a.s. for all $n$ large enough therefore $\P(X_\infty(k) = \infty) = 0$. Alternatively, if
\[
\sum_{n \geq 1} \P(B_n) = \infty,
\]
because the events are negatively correlated, we have that $B_n$ holds for infinitely many $n$ a.s. (c.f. \cite{KochenStone64}), which shows that $\P(X_\infty(k) = \infty) = 1$. Therefore, we just need to prove that, for all $n < m$, we have
\[
\P(B_n \cap B_m) \leq \P(B_n)\P(B_m).
\]
Note that up to a shift in the initial condition, we can assume without loss of generality that $k=0$. Additionally, using the Markov property, we have
\[
\P(B_n \cap B_m\mid X_1,\ldots, X_{n-1}) = \P_{X_{n-1}}(B_{m-n+1} \cap B_1) = \P_{X_{n-1}}(B_1) \P_{X_{n-1}}(B_{m-n+1}\mid B_1)
\]
with the usual notation $\P_X$ to denote a probability under which the IBM starts from configuration $X$. Therefore, it is enough to show that for any initial configuration, and all $m \geq 1$, we have
\begin{equation}\label{eq:01negcor}
\P_{X}(B_{m+1}\mid B_1) \leq \P_{X}(B_{m+1}).
\end{equation}
We write $A, B=B_1, C$ the events where the first ball is added to a bin of negative, null, or positive index respectively. We first check that
\begin{equation}\label{eq:condiBC1}
\P_{X}(B_{m+1}\mid B_1) = \P_{X}(B_{m+1}\mid C).
\end{equation}
This equality follows from Lemma~\ref{lemma:couplage1} with $j_0
:= -1$, $X_0$ the configuration obtained from $X$ by adding a ball in a bin at a positive index and $\tilde{X}_0$ the configuration obtained from $X$ by adding a ball in bin $0$. With this setting, \eqref{eq:condiBC1} is a direct consequence of equivalence~\eqref{incH1} applied at $j = j_0 + 1 = 0$.

We also claim that
\begin{equation}\label{eq:condiBC2}
\P_X(B_{m+1}\mid A) \geq \P_X(B_{m+1}\mid B_1).
\end{equation}
This inequality now follows from Lemma~\ref{lemma:couplage2}, with $j_0
:= -1$, $X_0$ the configuration obtained from $X$ by adding a ball in a bin of negative index and $\tilde{X}_0$ the configuration obtained from $X$ by adding a ball in bin $0$. With this setting, \eqref{eq:condiBC2} is now a direct consequence of implication~\eqref{incH2} applied again at $j = j_0 + 1 = 0$.

Finally, combining~\eqref{eq:condiBC1} and~\eqref{eq:condiBC2} we find that
\begin{align*}
\P_X(B_{m+1}) & = \P_X(B_{m+1}\mid A)\, \P_X(A) + \P_X(B_{m+1}\mid B)\, \P_X(B)\\
 &\qquad + \P_X(B_{m+1}\mid C)\, \P_X(C) \\
& \geq \P_X(B_{m+1}\mid B) \big(\P_X(A) + \P_X(B) + \P_X(C) \big)\\
& = \P_X(B_{m+1}\mid B)
\end{align*}
which establishes~\eqref{eq:01negcor} hence completes the proof of the of Theorem~\ref{thm:01law}.
\end{proof}
\begin{thebibliography}{ČS24}
\bibitem[FK03]{FK} Sergey Foss and Takis Konstantopoulos, Extended renovation theory and limit theorems for stochastic ordered graphs, Markov Process. Relat. Fields 9 (2003), no. 3, 413–468.
\bibitem[KS64]{KochenStone64} Simon Kochen and Charles Stone, A note on the Borel–Cantelli lemma, Ill. J. Math. 8 (1964), 248–251.
\end{thebibliography}
\end{document}
